\chapter{Fundamentação e trabalhos relacionados}
\label{cap:fundamentacao}

A avaliação proposta envolve três níveis: a relação que se deseja verificar, sua representação em circuitos e o custo de executar os componentes do sistema. Separar esses níveis permite interpretar por que uma mudança de tipo pode modificar a implementação sem mudar a relação, enquanto a adição de um compromisso modifica a própria afirmação verificada.

\section{Provas de conhecimento zero e provas de intervalo}

Uma relação $R(x,w)$ associa uma instância pública $x$ a uma testemunha $w$. O provador utiliza a testemunha para sustentar a afirmação, enquanto o verificador avalia sua aceitação. A referência da comunidade ZKProof distingue completude, correção contra provadores desonestos e conhecimento zero. A completude trata da aceitação quando a afirmação e a execução honesta são válidas; a correção limita a aceitação de afirmações falsas. Conhecimento zero é definido por simulação e limita a informação adicional revelada sobre a testemunha. Essas propriedades possuem definições formais e dependem do modelo adotado \cite{zkproofref}.

Neste trabalho, a instância pública inclui os limites do intervalo, e a testemunha inclui o valor privado. Uma prova de intervalo estabelece que esse valor pertence a um intervalo especificado. A proposta de benchmarking da comunidade ZKProof inclui essa funcionalidade e distingue a relação de alto nível da representação utilizada para implementá-la \cite{zkbenchmark}. A distinção é relevante porque representar um valor em um corpo finito não impõe, por si só, a interpretação inteira pretendida pelo intervalo.

A expressão ``prova'' será utilizada para os artefatos gerados pelo sistema estudado. A aceitação dos testes de implementação não será apresentada como demonstração das propriedades formais do protocolo. Da mesma forma, manter um valor como entrada privada não torna privados os limites, identificadores ou demais dados explicitamente públicos.

\section{Compromissos e relações verificadas}

Para o protótipo, um compromisso público permite formular uma relação que liga o valor privado a um objeto já conhecido pelo verificador. A afirmação deixa de ser apenas ``existe um valor no intervalo'' e passa a exigir que o valor usado na comparação também participe do cálculo do compromisso fornecido. Essa ligação é expressa no próprio circuito, conforme a Equação~\ref{eq:compromisso}.

A implementação utiliza uma função Pedersen da biblioteca adotada, um valor auxiliar privado e um identificador público. A formulação é específica do protótipo e não deve ser confundida automaticamente com qualquer construção denominada compromisso de Pedersen na literatura. Sua finalidade experimental é medir a relação adicional implementada. Autenticação de emissor, origem do atributo e proteção contra reutilização exigiriam mecanismos além dessa igualdade.

As provas de intervalo sobre valores comprometidos são um tema da literatura criptográfica. Bulletproofs apresenta uma construção não interativa sem configuração inicial confiável, com tamanho de prova logarítmico na dimensão considerada, e trata explicitamente de intervalos sobre valores comprometidos \cite{bulletproofs}. Essa construção fornece contexto para a funcionalidade, mas não é o protocolo executado neste estudo.

\section{Noir, ACIR, UltraHonk e ambientes de execução}

Noir é uma linguagem específica de domínio voltada à construção de programas que utilizam provas de conhecimento zero \cite{noirdocs}. No fluxo deste projeto, o programa é compilado para ACIR e posteriormente processado pelo backend Barretenberg para geração das provas UltraHonk. A cadeia efetivamente utilizada, seus artefatos e versões são apresentados na Seção~\ref{cap:implementacao}.

A documentação pública do Noir consultada identifica uma versão posterior à instalada neste experimento. Ela é utilizada aqui apenas para caracterizar a finalidade da linguagem. As afirmações sobre APIs, relações e configuração executada dependem dos artefatos e fontes da versão congelada. Não se presume que uma API atual descreva retroativamente o comportamento da versão experimental.

Uma representação intermediária e um backend cumprem funções distintas. Benarroch et al. descrevem a transformação de programas em representações verificáveis e o processamento posterior pelo sistema de provas, ressaltando que os custos dependem das duas camadas \cite{zkbenchmark}. Essa separação orienta a apresentação conjunta de operações ACIR e gates: as duas contagens não são unidades intercambiáveis, e nenhuma delas determina isoladamente a duração observada.

Na documentação associada à versão adotada do Barretenberg, Honk combina restrições polinomiais do circuito, o protocolo Sumcheck e compromissos polinomiais KZG. Sumcheck reduz uma afirmação sobre uma soma de avaliações a verificações de avaliações polinomiais. A etapa de compromissos emprega Gemini para transformar afirmações multilineares em univariadas e Shplonk para agrupá-las antes da abertura KZG. A configuração Ultra inclui relações aritméticas, de permutação, de consulta a tabelas e de intervalo, entre outras \cite{ultrahonk}.

Essa arquitetura distingue a relação escrita pelo programador das operações executadas pelo provador. O transcript serializa mensagens e deriva desafios pela transformação Fiat--Shamir; as configurações destinadas à verificação EVM utilizam Keccak. As variantes com conhecimento zero incorporam mecanismos de mascaramento \cite{ultrahonk}. No experimento, o alvo EVM é mantido nos dois ambientes, conforme o método, para preservar a configuração criptográfica da comparação.

WebAssembly é um formato de código de baixo nível projetado para portabilidade e execução eficiente, utilizável em browsers e em outros ambientes. A especificação distingue o núcleo da linguagem das interfaces fornecidas pelo ambiente que a hospeda \cite{wasmspec}. Seus objetivos de desempenho não constituem garantia de igualdade com uma execução nativa para qualquer aplicação. Neste estudo, a comparação envolve um processo de linha de comando e um runtime residente no browser, com diferenças de carregamento e instrumentação descritas no método.

\section{Verificação EVM e componentes de custo}

A utilização de verificadores Solidity introduz restrições distintas do custo de geração da prova. A EIP-170 limita a 24.576 bytes o código retornado pela inicialização do contrato, no regime em que essa regra se aplica \cite{eip170}. Portanto, tamanho do código implantado e gas de verificação são medidas complementares: atender ao limite de código não demonstra baixo custo de execução.

A composição da calldata também interfere no custo de uma transação. A EIP-7623 define um piso dependente dos bytes de dados, combinando o componente convencional com um mínimo por tokens de calldata \cite{eip7623}. Quando esse piso determina o total, o recibo não identifica univocamente o componente de execução por simples subtração. A aplicação da regra ao experimento depende da confirmação do hardfork efetivamente executado, não apenas de consultar a fórmula.

A EIP-4844 introduz transações com blobs e uma contabilização específica. Os dados completos dos blobs não ficam diretamente acessíveis à execução EVM como calldata; a especificação utiliza compromissos e interfaces próprias \cite{eip4844}. Por isso, projetar ocupação de bytes em blobs não equivale a implementar uma substituição direta dos dados consumidos pelo verificador nem demonstra economia monetária.

\section{Trabalhos relacionados e posicionamento}

Os trabalhos selecionados cumprem funções diferentes nesta pesquisa. Bulletproofs trata da construção criptográfica e de provas de intervalo sobre compromissos, incluindo agregação \cite{bulletproofs}. O presente estudo utiliza circuitos Noir e UltraHonk, sem implementar Bulletproofs como concorrente. Consequentemente, seus resultados publicados não serão colocados ao lado das medições locais como se constituíssem uma comparação controlada de desempenho.

A proposta de Benarroch et al. organiza desafios de comparação de sistemas de provas e múltiplas métricas de custo \cite{zkbenchmark}. O recorte deste TCC é mais restrito: manter a cadeia de ferramentas e comparar relações e modos de execução específicos. A contribuição pretendida é uma avaliação rastreável dessa configuração, e não um ranking geral de sistemas criptográficos. A proposta comunitária é uma referência metodológica, não uma certificação do procedimento local.

Kalibera e Jones estudam fontes de variação em benchmarking de sistemas e a distribuição de repetições entre níveis como construções, execuções e iterações \cite{kalibera}. Essa perspectiva fundamenta a necessidade de distinguir ciclos internos de sessões reiniciadas. Entretanto, o artigo não prescreve os 20 dias deste experimento nem valida automaticamente o estimador de mediana diária escolhido. A alocação e as premissas específicas são justificadas na Seção~\ref{cap:metodo}.

No eixo de abrangência da avaliação, zk-Bench propõe examinar diferentes camadas da implementação, desde bibliotecas aritméticas até ferramentas de desenvolvimento de SNARKs, e inclui configurações distintas de hardware \cite{zkbench}. Esse escopo contrasta com o desenho deste TCC, que fixa uma máquina e uma cadeia de ferramentas. A comparação local concentra-se nas relações implementadas e nos ambientes de execução; a cobertura de oito variantes não representa diversidade de plataformas criptográficas.

No eixo browser/nativo, Jangda et al. avaliam aplicações da suíte SPEC CPU com Browsix-Wasm e encontram diferenças de desempenho entre código nativo e WebAssembly executado em browsers \cite{jangda}. O trabalho oferece um antecedente para a avaliação empírica desse eixo, mas suas cargas não são os provadores deste estudo. Versões dos motores, aplicações e configuração de execução delimitam os resultados publicados, impedindo que seus fatores de desempenho sejam usados como previsão para UltraHonk.

Há também avaliações diretamente voltadas a provadores Noir. O repositório ZKEmail.nr Benchmarks distingue navegador/WASM, Node.js e execução nativa para um circuito de verificação DKIM, apresentando separadamente primeira execução, agregados posteriores e geração de testemunha \cite{zkemailbench}. Essa separação evita identificar toda execução JavaScript como execução em navegador. Contudo, no código web consultado, o acumulador exclui a primeira execução e soma uma execução a menos que o divisor usado para a média. A correspondência desse código com os valores publicados não está demonstrada; os agregados não são utilizados aqui para estimar fatores de desempenho.

Duff compara provas de pertencimento e não pertencimento em árvores Merkle nos modos nativo e WASM, declarando entradas e artefatos compilados idênticos. A metodologia separa geração de testemunha, geração de prova e inicialização, excluindo esta última das tabelas de dez execuções \cite{duff}. Os trechos metodológicos e tabelas consultados não identificam suficientemente hardware, navegador e threads para reconstruir o ambiente da comparação. A contribuição para este TCC é a explicitação de fases e controles declarados; a interpretação local permanece vinculada às fronteiras de medição e à configuração efetivamente registradas neste experimento.

No eixo geração/verificação, o relatório técnico de Nadolinski para a Base compara SnarkJS, Rapidsnark, Gnark e Noir/Barretenberg na prova de verificação de assinaturas ECDSA P-256, medindo tempo de geração e gas em instâncias AWS EC2 \cite{basebenchmark}. O relatório identifica a execução de Noir usada nos benchmarks como nativa, embora mencione suporte a browser. A proximidade com este TCC está nas métricas e na presença de UltraHonk; a relação criptográfica e o hardware são diferentes. Por isso, o relatório contextualiza os componentes de custo, sem fornecer uma comparação controlada com nossas provas de intervalo.

No plano da documentação experimental, a campanha V1 do ProveKit define uma fronteira entre entrada estruturada e prova serializada, distingue runtime novo de runtime reutilizado e mantém indisponibilidades separadas de valores medidos \cite{provekitcampaign}. A superfície publicada utiliza Chrome/WASM em um Mac e execução nativa em dispositivos móveis; portanto, esses alvos não isolam o efeito do ambiente no mesmo hardware. O documento também explicita diferenças entre as afirmações provadas pelos sistemas comparados. Sua utilidade para este estudo está na rastreabilidade das condições e na delimitação das comparações, não na transferência de resultados numéricos.

No plano da aplicação, Gasser, Humbert e Elghareeb apresentam uma prova de conceito de credenciais para identidade eletrônica que inclui um predicado de idade em Noir \cite{eidreport}. O relatório distingue esse predicado dos mecanismos de vínculo com o titular e de verificação da assinatura do emissor, indicando que precisam ser combinados para sustentar a apresentação de uma credencial válida. Essa distinção ajuda a delimitar o protótipo deste TCC: verificar um intervalo e sua consistência com um compromisso não implementa essas demais garantias. Os custos do predicado isolado não são adotados como custos de uma credencial completa.

A seleção reúne três eixos complementares: camadas do sistema, ambiente de execução e custo de geração/verificação. O recorte deste trabalho combina esses eixos em uma configuração delimitada, com relações explicitamente descritas e medições rastreáveis. A busca realizada é dirigida, não sistemática; a seleção não sustenta alegação de exclusividade ou ausência de estudos semelhantes.
